\clearpage
\SourceLeaf{76}
\Needspace{10\baselineskip}\begin{center}
\textbf{[VIER NOTIZEN ÜBER INVERSION DER POTENZREIHEN.]}
\end{center}



\Needspace{10\baselineskip}\begin{center}
\textsc{Gauss} an \textsc{Olbers}, 12.\ December 1813.
\end{center}

\noindent\hspace*{\parindent}.\,.\,.\,.\,.\,.\ Heute habe ich einen artigen kleinen Fund gemacht, den ich Ihnen doch mittheilen
will. Es ist ein specieller sehr kurzer Beweis für die Umkehrungsformel der Reihen: man kann sie zwar
als einen einzelnen Fall sehr leicht aus der allgemeinen \textsc{Lagrang}ischen Umkehrungsformel
ableiten; aber es ist doch angenehm, sie mit wenigen Federstrichen unabhängig von dieser ableiten zu
können.

\noindent\hspace*{\parindent}Ich setze
\[
\begin{array}{l}
x+\alpha xx+\beta x^3+\gamma x^4+\text{etc.}=X,\\[4pt]
y+\alpha yy+\beta y^3+\gamma y^4+\text{etc.}=Y,\\[4pt]
\log\Bigl(1-\dfrac{x}{y}\Bigr)-\log\Bigl(1-\dfrac{X}{Y}\Bigr)=\Omega,\\[4pt]
Y^{-n}=y^{-n}\bigl(1+(n,1)y+(n,2)yy+(n,3)y^3+\text{etc.}\bigr).
\end{array}
\]

\noindent Man hat erstlich:
\[
\begin{array}{l}
\Omega=\dfrac{X}{Y}+\dfrac{X^2}{2Y^2}+\dfrac{X^3}{3Y^3}+\text{etc.}\\[6pt]
\hphantom{\Omega=}-\dfrac{x}{y}-\dfrac{x^2}{2y^2}-\dfrac{x^3}{3y^3}-\text{etc.}\\[6pt]
=\dfrac{X}{y}+(1,1)X+(1,2)Xy+.\,.\,.\,.\\[6pt]
\hphantom{=}+\dfrac{X^2}{2yy}+(2,1)\dfrac{X^2}{2y}+(2,2)\dfrac{X^2}{2}+(2,3)\dfrac{X^2y}{2}+.\,.\,.\,.\\[6pt]
\hphantom{=}+\dfrac{X^3}{3y^3}+(3,1)\dfrac{X^3}{3yy}+(3,2)\dfrac{X^3}{3y}+(3,3)\dfrac{X^3}{3}+(3,4)\dfrac{X^3y}{3}+.\,.\,.\,.\\[6pt]
.\,.\,.\,.\,.\,.\,.\,.\,.\,.\,.\,.\,.\,.\,.\,.\,.\,.\,.\,.\,.\,.\,.\,.\\[6pt]
.\,.\,.\,.\,.-\dfrac{x^3}{3y^3}-\dfrac{xx}{2yy}-\dfrac{x}{y}.
\end{array}
\]

